%!TEX program = uplatex
\RequirePackage{plautopatch}
\documentclass[uplatex,dvipdfmx]{jlreq}
\usepackage{amsmath,amssymb,amscd,amsthm}
\usepackage{url} % or hyperref
\pagestyle{empty}

\begin{document}

%======================================================================
% 表紙
%======================================================================
\begin{center}
  \vspace*{20pt}
  {\Huge\bfseries Poincaré 予想と３次元トポロジー}

  \vspace{1pc}
  {\Large 概要}
\end{center}

\vfill

\newpage

%======================================================================
% 講演１：小島定吉「ポアンカレ予想から幾何化予想へ」
%======================================================================
\begin{center}
  {\Large\bfseries ポアンカレ予想から幾何化予想へ}
\end{center}

\bigskip
\hspace*{\fill}{\large\bfseries 小島　定吉（東工大・情報理工）}

\bigskip

ポアンカレ予想とは、
約100年前に多様体のトポロジーの研究の基礎を確立した
ポアンカレの一連の大論文の中に記された
「単連結３次元閉多様体は３次元球面と位相同型か？」
という問題を指す。
いまだ未解決であるが、
20世紀の数学研究推進の大きな原動力となった予想で、
その重要性は衆目の一致するところである。

今回の ENCOUNTER with MATHEMATICS は、
ポアンカレ予想をめぐるこれまでの研究の進展と、
３次元という今注目が集まっている次元でのトポロジーの現況をテーマに取り上げている。
この講演は表題のとおり、
そもそものポアンカレ予想を出発点として、
３次元多様体論の進展を時代を追って解説し、全体の序章としたい。

\vspace{1pc}

\begin{center}
  \textsc{参考文献}
\end{center}
\begin{enumerate}
  \item[{[1]}] 加藤十吉，「ポアンカレ予想周辺」，数学，岩波書店，\textbf{31} (1979).
  \item[{[2]}] 小島定吉，「ポアンカレ予想」，数学セミナー，4月号 (1998).
  \item[{[3]}] 松本幸夫，「４次元ポアンカレ予想の解決」，数学セミナー，6月号 (1982).
  \item[{[4]}] J.~Milnor, The Poincaré Conjecture,\\
               \hspace*{\fill}\url{http://www.claymath.org/prize_problems/index.htm}
\end{enumerate}

\newpage

%======================================================================
% 講演２：加藤十吉「高次元ポワンカレ予想の解決」
%======================================================================
\begin{center}
  {\Large\bfseries 高次元ポワンカレ予想の解決}
\end{center}

\bigskip
\hspace*{\fill}{\large\bfseries 加藤　十吉（九州大学・数理学研究院）}

\bigskip

ポワンカレ予想の高次元版
「$n$（$\geq 5$）次元のホモトピー球面は $n$ 次元球面に同相である」
が S.~Smale によって解決されたのは周知の事実だろう。
しかし、文献的には、J.~Stallings の証明の方が早く出版されている。
そこには高次元ポワンカレ予想の解決を巡る
組合せ位相幾何学対微分位相幾何学の凄まじい先陣争いが伺える。
Smale は英国の一匹狼 J.~H.~C.~Whitehead の「collapsing と正則近傍の理論」
と米国の M.~Morse の理論および H.~Whitney の「可微分埋蔵の交叉理論」
の融合により、高次元ポワンカレ予想に到達したと言われる。
Smale の理論はそれだけの解決にとどまることなく、一般の高次元
多様体の分類理論の基盤となる $h$-同境定理に到達する。

\vspace{1pc}

今回の ENCOUNTER with MATHEMATICS では、Smale 以前の
Whitehead の高次元ポワンカレ予想への取り組みを中心に、
組合せ位相幾何的な、したがって、
ユークリッドの原本の第13巻としての高次元ポワンカレ予想の解決に触れて、
Stallings の組合せ位相幾何学的証明への執着と情熱を理解する一助として
戴ければ幸です。

\newpage

%======================================================================
% 講演３：松本幸夫「４次元ポアンカレ予想の解決」
%======================================================================
\begin{center}
  {\Large\bfseries ４次元ポアンカレ予想の解決}
\end{center}

\bigskip
\hspace*{\fill}{\large\bfseries 松本　幸夫（東大・数理）}

\bigskip

1904年にポアンカレによって３次元ポアンカレ予想が問題とされて以来、
約60年後にスメールにより高次元の類似（高次元ポアンカレ予想）が解決された。
これは、「ある意味では高次元のほうが見やすい」という驚くべき事実の発見だった。
スメールの方法は高次元に特有の技法、とくに「ホイットニー・トリック」を使うため、
すぐに４次元に適用するわけには行かなかったが、
1973年のキャッソンの素晴らしいアイデアによって、
４次元に適用する道が開かれた。
そして遂に1981年になり、フリードマンによる４次元ポアンカレ予想の解決をみた。

この講演では、キャッソンのアイデアとフリードマンの仕事についての解説を試みたいと思います。

\vspace{1pc}

\begin{center}
  \textsc{参考文献}
\end{center}
\begin{enumerate}
  \item[{[1]}] 上正明・久我健一，
        「Freedman による４次元 Poincaré 予想の解決について」，
        数学 \textbf{35}（1983）.
\end{enumerate}

\newpage

%======================================================================
% 講演４：小島定吉「３次元トポロジーから三題 I. 幾何化予想から ?? へ」
%======================================================================
\begin{center}
  {\Large\bfseries ３次元トポロジーから三題}\\[2pc]
  {\large\bfseries I.\quad 幾何化予想から ?? へ}
\end{center}

\bigskip
\hspace*{\fill}{\large\bfseries 小島　定吉（東工大・情報理工）}

\bigskip

幾何化予想は、
ポアンカレ予想をいっきに３次元多様体のトポロジーの分類に引き上げる
新しい視点をもたらした。
こうした幾何的な立場からの研究は、
基本群が無限群である場合に著しく進展し、
近い将来の解決も夢ではないと期待されている。
しかし基本群が有限群の場合はどうも相性が今一つという感が否めない。

一方こうした幾何化の流れとは独立に、
80年代半ばから新しい代数的手法による
３次元多様体の位相不変量の研究が
数理物理の興味と結びつき大きく展開した。
それぞれの研究は長い間（超）平行線を歩んだが、
最近になって、体積予想など両者を結びつける橋が現われ始めてきた。
この講演では、
これからどこへ向かうのか見当もつかない３次元トポロジーを、
こうした橋を意識しつつ（すでに古典かもしれない）幾何化の
視点から今一度振り返ってみたい。

\newpage

%======================================================================
% 講演５：大槻知忠「３次元トポロジーから三題 II. 3次元多様体の量子不変量」
%======================================================================
\begin{center}
  {\large\bfseries II.\quad 3次元多様体の量子不変量}
\end{center}

\bigskip
\hspace*{\fill}{\large\bfseries 大槻　知忠（東工大・情報理工）}

\bigskip

1980年代の数理物理と低次元トポロジーの交流により、
3次元トポロジーには、量子不変量とよばれる多数の位相不変量がもたらされた。
数理物理的には、Chern--Simons 汎関数を Lagrangian とする経路積分として
量子不変量は表される。
この Chern--Simons 経路積分を数学的に理解する一つの方法が摂動展開であり、
この摂動展開を数学的に定式化することにより
3価グラフの無限線形和が（3次元多様体の位相不変量として）えられる。
とくにホモロジー球面に話を限ると、
この無限線形和の各係数は整数性をもつ有限型不変量になり、
さらに各有限型不変量は任意の整数値をとる。
また、すべての有限型不変量が同じであるような同相でないホモロジー球面は
現時点では知られていない。
希望的には、ホモロジー球面の集合は
3価グラフがはる（無限 rank の）lattice の部分集合とみなされる（ことが期待される）。
量子不変量はこの lattice 上の関数である。

最近の話題である体積予想によると、量子不変量のある種の極限として
（数理物理的にはある種の相関関数の準古典極限として）
「双曲体積」と「Chern--Simons 不変量」が現われる（らしい）ことが
数値実験等により強く示唆されている。
今後、幾何構造と位相不変量の研究が結びつくことにより、
3次元多様体のトポロジーの解明にむけて、
体積予想から何がもたらされるのか？ 今後の進展が期待される。

この講演では、Chern--Simons 経路積分から体積予想まで、
数理物理に由来する不変量をめぐる3次元トポロジーの
（この十数年の）発展を概観する。

\newpage

%======================================================================
% 講演６：吉田朋好「３次元トポロジーから三題 III. ３次元多様体と２次特性類」
%======================================================================
\begin{center}
  {\large\bfseries III.\quad ２次元から見た量子不変量}
\end{center}

\bigskip
\hspace*{\fill}{\large\bfseries 吉田　朋好（東工大・理）}

\bigskip

3次元量子不変量を2次元共形場理論から見たときの風景について話す。

\end{document}
